\documentclass[11pt]{article}
\usepackage{amsmath,amssymb,booktabs,geometry}
\geometry{margin=1in}

\title{Step 197: Mertens Criterion Carrier Pivot}
\author{six-birds-foundations-iii}
\date{}

\begin{document}
\maketitle

\section*{Verdict}

\[
\boxed{\texttt{V\_mertens\_CRCFT\_TE}.}
\]

The weak Mertens criterion is a classically-RH-equivalent carrier and lands in
CRCFT target-equivalence mode.  The strong Mertens conjecture is a separate,
strictly stronger assertion and is classically refuted.

\section*{Inherited Record Audit}

No direct Mertens carrier record was found in the inherited RH thread,
\texttt{anti\_loc/adequacy.tex}, or \texttt{anti\_loc/needles.tex}.  The closest
inherited object is Step 131's finite Boolean Mobius-annihilation identity:
\[
  \langle\psi_-,Z_yb\rangle=2^{-k/2}(-1)^k b(P_y),
\]
which is an exact finite incidence identity.  It is not an asymptotic Mertens
carrier and does not supply bounds for \(M(x)\).

\section*{Carrier Declaration}

Let \(\mu(n)\) be the Mobius function and define
\[
  M(x)=\sum_{n\le x}\mu(n).
\]
The typed carrier is the discrete/asymptotic sequence ledger generated by
\(\mu\), the partial-sum operator \(S\mu=M\), and asymptotic envelopes for the
readout \(M(x)\).  Native probes include:
\[
  \mu(n),\qquad
  \sum_{n\in I}\mu(n),\qquad
  \sum_{\substack{n\le x\\ n\equiv a\pmod q}}\mu(n),
\]
and the Dirichlet-series probe
\[
  \frac1{\zeta(s)}=\sum_{n=1}^{\infty}\frac{\mu(n)}{n^s}
  \qquad(\Re s>1).
\]
The dissolving probes are the barrier functions \(x^{1/2+\epsilon}\), for
\(\epsilon>0\).

\section*{Residual}

For \(\epsilon>0\), define
\[
  \Xi_M(\epsilon)=
  \limsup_{x\to\infty}\frac{|M(x)|}{x^{1/2+\epsilon}}.
\]
Closure of the Mertens carrier means
\[
  \Xi_M(\epsilon)=0\quad\text{for every }\epsilon>0.
\]
Equivalently, \(M(x)=O(x^{1/2+\epsilon})\) for every \(\epsilon>0\).

\section*{Strong versus Weak Mertens}

\begin{center}
\begin{tabular}{lll}
\toprule
Assertion & Status & Framework role\\
\midrule
\(|M(x)|<\sqrt{x}\) & refuted & strong sidecar no-go\\
\(M(x)=O(x^{1/2+\epsilon})\ \forall\epsilon>0\) & open / RH-equivalent & primary CRE residual\\
\bottomrule
\end{tabular}
\end{center}

Mertens' 1897 conjecture asserted the square-root bound.  Odlyzko--te Riele
(1985) refuted this strong form by proving that the normalized summatory
Mobius ratio crosses the square-root barrier in magnitude.  This does not
refute RH, because the strong form is strictly stronger than RH.

Littlewood's 1912 criterion gives the weak equivalence:
\[
  \mathrm{RH}
  \quad\Longleftrightarrow\quad
  M(x)=O(x^{1/2+\epsilon})\quad\text{for every }\epsilon>0.
\]
The usual route uses partial summation and the analytic behavior of
\(1/\zeta(s)\).  A bound of this strength gives the required continuation and
zero-free condition past \(\Re s>1/2\), while RH supplies the corresponding
summatory estimate with arbitrary \(\epsilon\).

\section*{CRCFT Audit}

The carrier is CRE because closure of \(\Xi_M\) is equivalent to Riemann RH.
Therefore the Dichotomy sends the carrier to CRCFT.  The primary mode is
\[
  \mathrm{CRCFT\mbox{-}TE}.
\]
This is a pre-reduction target-equivalence obstruction: closing the residual
is exactly the target theorem.

The CTMT mode is not primary.  Finite values of \(M(x)\) are computable and
can form a finite-data interface, but the missing step is asymptotic
promotion to every \(\epsilon>0\), not a carrier-native matrix element whose
kernel is absent.  The BF mode is also not primary: the Mertens carrier is
native and well-defined.

\section*{Dichotomy Implication}

The Mertens carrier adds another CRE instance supporting the Dichotomy and
CRCFT coverage pattern:
\[
  \text{CRE carrier}\quad\Rightarrow\quad
  \text{CRCFT foreclosure mode}.
\]
Here the foreclosure mode is target-equivalence, with an additional classical
sidecar: the stronger square-root conjecture has already been refuted.

\section*{Classical Sources}

\begin{itemize}
\item Stieltjes, 1885 correspondence/claims on Mobius summatory bounds.
\item Mertens, 1897, \emph{Uber eine zahlentheoretische Funktion}.
\item Littlewood, 1912, \emph{The converse of Abel's theorem on power series}.
\item Odlyzko and te Riele, 1985, \emph{Disproof of the Mertens conjecture}.
\end{itemize}

\section*{Nonclaims}

This step does not prove RH, disprove RH, prove the weak Mertens criterion, or
transfer Mertens information to the Burnol/Sonine carrier.  It classifies the
Mertens criterion carrier inside the existing CRCFT/Dichotomy framework.

\end{document}

